% 贯穿四点图 E={12,13,23,24,34}。S={2,3}：G-S 的分量 {1}、{4}；局部块 G[{1,2,3}]、G[{2,3,4}]，共享接口 {2,3} 与边 {2,3}。
% 对照：S={1} 时 G-S 是三角形 {2,3,4}，只有一个分量。
\begin{tikzpicture}[font=\figurefont\fontsize{8}{10}\selectfont,
  vtx/.style={circle,draw=FigureStroke,line width=.8pt,minimum size=6.2mm,inner sep=0pt,text=FigureInk,fill=FigureBlueFill},
  sep/.style={vtx,fill=FigureYellowFill},
  gone/.style={vtx,draw=FigureMuted,fill=white,text=FigureMuted,dashed},
  edge/.style={draw=FigureStroke,line width=.85pt},
  cut/.style={draw=FigureRed,line width=.8pt,dashed},
  link/.style={draw=FigureMuted,line width=.5pt,dotted},
  note/.style={font=\FigureNoteFont\fontsize{7.5}{9}\selectfont,text=FigureMuted,align=center}]
% (a)
\begin{scope}
  \node[vtx] (a1) at (0,0) {1}; \node[sep] (a2) at (1.0,0.75) {2};
  \node[sep] (a3) at (1.0,-0.75) {3}; \node[vtx] (a4) at (2.0,0) {4};
  \draw[cut] (a1)--(a2) (a1)--(a3) (a2)--(a3) (a2)--(a4) (a3)--(a4);
  \node[note] at (1.0,-1.45) {删除 $S=\{2,3\}$ 后\\剩下分量 $\{1\}$、$\{4\}$};
  \node at (1.0,-2.3) {(a) 删除候选集};
\end{scope}
% (b)
\begin{scope}[xshift=3.2cm]
  \node[vtx] (b1) at (0,0) {1}; \node[sep] (b2) at (0.95,0.75) {2}; \node[sep] (b3) at (0.95,-0.75) {3};
  \draw[edge] (b1)--(b2) (b1)--(b3) (b2)--(b3);
  \node[sep] (c2) at (2.05,0.75) {2}; \node[sep] (c3) at (2.05,-0.75) {3}; \node[vtx] (c4) at (3.0,0) {4};
  \draw[edge] (c2)--(c3) (c2)--(c4) (c3)--(c4);
  \draw[link] (b2)--(c2) (b3)--(c3);
  \node[note] at (1.5,-1.45) {$\mathcal G[\{1,2,3\}]$ 与 $\mathcal G[\{2,3,4\}]$\\点线连接同一原顶点};
  \node at (1.5,-2.3) {(b) 加回接口得到局部块};
\end{scope}
% (c)
\begin{scope}[xshift=7.6cm]
  \node[sep] (d1) at (0,0) {1}; \node[vtx] (d2) at (1.0,0.75) {2};
  \node[vtx] (d3) at (1.0,-0.75) {3}; \node[vtx] (d4) at (2.0,0) {4};
  \draw[cut] (d1)--(d2) (d1)--(d3);
  \draw[edge] (d2)--(d3) (d2)--(d4) (d3)--(d4);
  \node[note] at (1.0,-1.45) {删除 $\{1\}$ 后\\只剩一个分量};
  \node at (1.0,-2.3) {(c) 无效候选集};
\end{scope}
\end{tikzpicture}
